\documentclass[11pt]{article}
% Shared preamble for all daily exercises
% Update this file to change settings (padding, parindent, etc.) for all exercises

\usepackage[margin=0.75in]{geometry}
\usepackage{amsmath}
\usepackage{amssymb}

% Custom settings
\setlength{\parindent}{0pt}
\setlength{\parskip}{0.2cm}

\title{Daily Exercise}
\author{Dhruman Gupta}
\date{27th March}

\begin{document}

% Common header for daily exercises
% This file can be included in other LaTeX documents

\maketitle

\noindent
\textbf{Course:} CS-3330, Theory of Computation

\vspace{0.2cm}

\noindent
\textbf{Question:}\noindent 
Prove a lower bound for sorting a sequence of integers in the Turing machine model.

\textbf{Answer:}

Suppose the input is a sequence of $n$ integers, each written in binary using exactly $b$ bits. Then the total input length is
\[
N = \Theta(nb).
\]

We will show that any Turing machine that sorts such an input must take time $\Omega(N)$, and hence $\Omega(nb)$.

Let $M$ be a Turing machine that sorts sequences of integers. For the sake of contradiction, assume that on some input $x$, the machine halts without ever reading some input bit position $i$.

Now construct another input $x'$ by changing only the bit in position $i$, and leaving all other bits the same. Since the integers are written in fixed-length binary, $x'$ is still a valid input. Also, by changing a bit inside one of the integers, we can ensure that the input sequence is different, and hence the correctly sorted output is different.

But the computation of $M$ on $x$ and $x'$ is identical. Indeed, the machine never reads position $i$, so at every step it sees exactly the same information on both inputs. Hence it makes the same transitions and produces the same output on both inputs.

This is a contradiction, since $x$ and $x'$ have different correct sorted outputs. Therefore, on every input, a correct sorting machine must read every input bit.

A Turing machine can read at most one tape cell in one step. So reading all $N$ input symbols requires at least $N$ steps. Hence the running time is at least
\[
\Omega(N) = \Omega(nb).
\]

Therefore, sorting a sequence of $n$ integers of $b$ bits each in the Turing machine model has lower bound $\Omega(nb)$.

In particular, if $b = \Theta(\log n)$, then the lower bound becomes
\[
\Omega(n \log n).
\]

\end{document}
